Mechanistic interpretability has matured considerably in the context of language models, and recent work has extended concept-based analysis to vision-language models. Yet a systematic application of these tools to Vision Transformers trained without any language supervision is still missing from the literature. This section identifies the specific gaps that motivate our work.

\subsection{The Absence of Systematic SAE Analyses for Pure Vision Models}

Sparse Autoencoders have proven effective at decomposing the internal representations of language models into monosemantic features. Cunningham et al. \cite{Cunningham2024} and Bricken et al. \cite{Bricken2023} both report that SAEs trained on transformer MLP activations or residual stream vectors recover sparse, human-interpretable directions that influence model behaviour in predictable ways. Similarly, the automated circuit discovery methodology of Conmy et al. \cite{Conmy2023} has identified specific subgraphs responsible for well-defined language tasks.

What these contributions have in common is that they target language models exclusively. The tokens processed by these models are words; the features that emerge from SAE decomposition correspond to linguistic concepts; and verifying interpretability amounts to reading the tokens that activate each feature. In Vision Transformers, none of these conveniences apply. The tokens are image patches, the activations encode low-level visual statistics alongside high-level semantics, and there is no direct way to read off a human-understandable label from a learned sparse direction. Whether SAEs recover meaningful structure in purely visual activation spaces is therefore an empirical question that has not been systematically addressed in the mechanistic interpretability literature.

\subsection{The Spatial Structure of ViTs Creates New Analytical Challenges}

Language transformers process sequences of tokens with a clear positional ordering. The causal structure of information flow is relatively straightforward to trace: earlier tokens influence later ones, and the residual stream at each position accumulates contributions from preceding layers and attention heads. In Vision Transformers, the situation is fundamentally different in at least two respects.

First, the tokens represent spatial patches of a two-dimensional image. Self-attention operates over all pairs of patches simultaneously, with no preferred direction of information flow. As Raghu et al. \cite{Raghu2021} show, ViTs aggregate global information across spatial locations very early in the network---a property that has no direct analogue in language models, where global context is built up gradually across layers and positions. This spatial promiscuity makes it harder to localise which attention heads are routing which pieces of visual information, and to which destination.

Second, the final prediction is derived exclusively from the \texttt{[CLS]} token, which receives contributions from all patch tokens through attention. Identifying the circuit responsible for a particular classification decision requires tracing how evidence from specific spatial regions propagates and gets compressed into this single token. The causal analysis tools developed for language tasks are in principle applicable here, but neither their assumptions nor their practical effectiveness in the visual setting have been validated.

\subsection{The Cross-Modal Labelling Problem}

A subtler but equally important gap concerns evaluation. When Gandelsman et al. \cite{Gandelsman2024} and Haque et al. \cite{Haque2026} identify visual concepts in CLIP or in medical VLMs, they can label those concepts in natural language because the model being analysed was trained on paired image-text data. The shared embedding space provides a free, built-in mechanism for translating visual directions into linguistic descriptions.

In a ViT trained purely on image classification labels---or self-supervisedly via DINO---no such space exists. A feature identified by an SAE is just a direction in a high-dimensional activation space: it activates strongly on certain image patches, but the model itself offers no vocabulary for describing what those patches have in common. Assigning a human-readable label to each feature requires an external grounding mechanism. This evaluation bottleneck is not acknowledged in most mechanistic interpretability work on ViTs, and it means that the interpretability claims made for SAE features in the visual domain are harder to substantiate than equivalent claims in the language domain.

\subsection{The Research Gap We Address}

These gaps converge on a single open question: can Sparse Autoencoders recover interpretable, monosemantic features from the internal activations of a purely visual Vision Transformer, and how can those features be evaluated rigorously without built-in language supervision?

Our project targets this question directly. We train SAEs on the MLP output activations of selected layers of a pre-trained ViT-B/16, and use CLIP as an \textit{external} cross-modal evaluator: given the images that maximally activate each discovered feature, CLIP produces a linguistic description by zero-shot matching against a broad textual vocabulary. This decouples the interpretability analysis from the backbone model, avoiding the need for a vision-language architecture while still grounding the features in human-understandable concepts. The two concrete gaps we address are the absence of systematic SAE analyses for purely visual models, and the lack of a validated evaluation methodology when no internal text-image alignment is available.
